\section{Observer interfaces, accessible quantities, and exact decoding}
\label{sec:interfaces}

A single observation map determines which states produce the same output. In dynamical observability, the available evidence is an output trajectory generated by specified dynamics and inputs \parencite{Kalman1963,HermannKrener1977,Sontag1998}.

\begin{definition}[Observer interface]
An observer interface on a world state space \(X\) is a pair
\[
I=(O,\pi),
\]
where \(O\) is an accessible observation space and \(\pi:X\to O\) is the observation map.
\end{definition}
\begin{definition}[Observational equivalence]
For \(x,y\in X\), define
\[
x\sim_I y
\quad\Longleftrightarrow\quad
\pi(x)=\pi(y).
\]
\end{definition}
\begin{proposition}
Observational equivalence is reflexive, symmetric, and transitive.
\end{proposition}

\begin{proof}
Reflexivity follows from \(\pi(x)=\pi(x)\). Symmetry follows because equality is symmetric. Transitivity follows because \(\pi(x)=\pi(y)\) and \(\pi(y)=\pi(z)\) imply \(\pi(x)=\pi(z)\).
\end{proof}
The equivalence class of \(x\) is the fiber
\[
[x]_I=\pi^{-1}(\pi(x)).
\]
The class contains exactly the source states that the interface represents by the same accessible value. Equivalence is therefore relative to \(I\). A second interface may separate states that the first interface merges.

\subsection{Quantities that factor through the interface}

An accessible quantity is a set function of the observation. Computability, cost, and stability under measurement error require additional hypotheses.

\begin{definition}[Interface accessible quantity]
Let \(q:X\to V\). The quantity \(q\) factors through \(I\) when there exists a map \(\bar q:O\to V\) such that
\[
q=\bar q\circ\pi.
\]
\end{definition}
\begin{theorem}[Fiber constancy]
If \(q\) factors through \(I\), then \(q\) is constant on every observation fiber. In particular,
\[
x\sim_I y
\quad\Longrightarrow\quad
q(x)=q(y).
\]
\end{theorem}

\begin{proof}
Assume \(q=\bar q\circ\pi\). Let \(x\sim_I y\). By definition, \(\pi(x)=\pi(y)\). Applying \(\bar q\) to both sides gives
\[
\bar q(\pi(x))=\bar q(\pi(y)).
\]
Substitution of \(q=\bar q\circ\pi\) yields \(q(x)=q(y)\).
\end{proof}
A quantity constant on every fiber has a unique representation on the observable quotient (\Cref{sec:observablequotient}).

\begin{corollary}[Variation prevents interface factorization]
If \(x\sim_I y\) and \(q(x)\neq q(y)\), then \(q\) does not factor through \(I\).
\end{corollary}

\begin{proof}
If \(q\) factored through \(I\), fiber constancy would imply \(q(x)=q(y)\), contradicting the assumption.
\end{proof}
\subsection{Internal tests}

An internal test is a special case of an accessible quantity.

\begin{definition}[Internal test]
An internal binary test is a function
\[
t:O\to\{\mathrm{false},\mathrm{true}\}.
\]
\end{definition}
\begin{proposition}[No internal test separates one fiber]
If \(x\sim_I y\), then
\[
t(\pi(x))=t(\pi(y))
\]
for every internal test \(t\).
\end{proposition}

\begin{proof}
Apply \(t\) to both sides of \(\pi(x)=\pi(y)\).
\end{proof}
Every test on \(O\) receives the same input for states in one fiber. Changing its output space cannot distinguish these states.

\subsection{Exact reconstruction}

An exact decoder is a left inverse of the observation map \parencite[Chs. 10--11]{Halmos1974}.

\begin{definition}[Globally exact decoder]
A globally exact decoder for \(I\) is a function \(d:O\to X\) such that
\[
d\circ\pi=\id_X.
\]
\end{definition}
\begin{theorem}[Exact decoding implies injectivity]
If a globally exact decoder exists, then \(\pi\) is injective.
\label{thm:decoderinjective}
\end{theorem}

\begin{proof}
Assume \(d\circ\pi=\id_X\). Let \(x,y\in X\) and assume \(\pi(x)=\pi(y)\). Applying \(d\) gives
\[
d(\pi(x))=d(\pi(y)).
\]
The left inverse condition reduces this equality to \(x=y\). Therefore \(\pi\) is injective.
\end{proof}
\begin{corollary}[Noninjectivity prevents exact decoding]
If \(\pi\) is noninjective, then no globally exact decoder exists.
\end{corollary}

\begin{proof}
The statement is the contrapositive of \Cref{thm:decoderinjective}.
\end{proof}
A many to one observation may still identify a particular target or identify a state within a restricted source class.

\subsection{Finite capacity obstruction}

For finite sets, counting the possible outputs can rule out injectivity.

\begin{theorem}[Finite pigeonhole obstruction]
Let \(X\) and \(O\) be finite. If
\[
|O|<|X|,
\]
then every observation map \(\pi:X\to O\) is noninjective. Therefore there exist distinct \(x,y\in X\) such that \(x\sim_I y\), and no globally exact decoder exists.
\end{theorem}

\begin{proof}
Assume for contradiction that \(\pi\) is injective. An injective map from a finite set \(X\) into a finite set \(O\) implies \(|X|\leq |O|\). This contradicts \(|O|<|X|\). Hence \(\pi\) is noninjective. The collision and decoder conclusions follow from the definitions and the preceding corollary.
\end{proof}

